% !TEX program = xelatex
% =============================================================================
%  ELSA Dataset — Multi-Word Entity Detection: Method and Comparative Analysis
% -----------------------------------------------------------------------------
%  MUST be compiled with XeLaTeX (Overleaf: Menu -> Compiler -> XeLaTeX).
%  Bengali font: Noto Sans Bengali (bundled on Overleaf / TeX Live).
%  On local Windows, replace with:  \newfontfamily\bnfont{Nirmala UI}
% =============================================================================
\documentclass[11pt]{article}

\usepackage[a4paper,margin=2.4cm]{geometry}
\usepackage{fontspec}
\newfontfamily\bnfont{Noto Sans Bengali}[Script=Bengali]
\newcommand{\bn}[1]{{\bnfont #1}}

\usepackage{xcolor}
\definecolor{brandblue}{HTML}{1F4E79}
\definecolor{brandteal}{HTML}{2E7D74}
\definecolor{codegreen}{HTML}{2E7D32}
\definecolor{codered}{HTML}{C62828}
\definecolor{noteblue}{HTML}{E7F0FA}

\usepackage[most]{tcolorbox}
\usepackage{booktabs}
\usepackage{array}
\usepackage{longtable}
\usepackage{enumitem}
\usepackage{titlesec}
\usepackage{hyperref}
\hypersetup{colorlinks=true, linkcolor=brandblue, urlcolor=brandteal}

\titleformat{\section}{\large\bfseries\color{brandblue}}{\thesection}{0.6em}{}
\titleformat{\subsection}{\normalsize\bfseries\color{brandteal}}{\thesubsection}{0.6em}{}

\newtcolorbox{note}[1][]{colback=noteblue,colframe=brandblue,boxrule=0.5pt,
  arc=2pt,left=8pt,right=8pt,top=6pt,bottom=6pt,fonttitle=\bfseries,title=#1}

\newcommand{\yes}{\textcolor{codegreen}{\textbf{Yes}}}
\newcommand{\no}{\textcolor{codered}{\textbf{No}}}

\setlist{leftmargin=1.4em,itemsep=2pt,topsep=3pt}

% =============================================================================
\begin{document}

\begin{center}
  {\LARGE\bfseries\color{brandblue} Multi-Word Entity Detection in the ELSA Dataset}\\[4pt]
  {\normalsize Detection method, results, and a comparison of automated vs.\ manual analysis}\\[3pt]
  {\small Bangla (Bengali) e-commerce product reviews\quad$\vert$\quad Pipeline: \texttt{elsa\_multiword\_entity\_bio\_v3.py}}\\[6pt]
  \rule{\textwidth}{0.8pt}
\end{center}

% =============================================================================
\section{Objective}

The goal was to identify \textbf{multi-word product entities} in 10,000 Bangla
product reviews --- product names composed of two or more tokens, such as
\bn{ডেটল সাবান} (Dettol Soap) or \bn{ফেস ওয়াশ} (Face Wash). Multi-word forms
matter because a lone token such as \bn{সাবান} (``soap'') is ambiguous, whereas
\bn{ডেটল সাবান} identifies both the brand and the product type.

Detection was performed two ways --- automatically with a Python pipeline and
manually by a human reviewer. The automated pipeline returned \textbf{13} unique
multi-word entities; the manual review returned \textbf{30}. This document
describes the detection method, reports the results, and explains the
difference.

% =============================================================================
\section{Data and Entity Tagging}

Each review carries a pre-tagged column, \texttt{ENTITY\_TAGGED\_REVIEW}, in
which every product token is marked with the prefix \texttt{\_NE\_} (Named
Entity). For example, in \bn{\_NE\_ডেটল} \bn{সাবান} \bn{তো} \bn{ভাল}, the token
\bn{ডেটল} carries the tag and begins a product name. The entire detection
pipeline depends on this tag: an untagged token is never treated as a product.

% =============================================================================
\section{Detection Method}

Detection is implemented in \texttt{elsa\_multiword\_entity\_bio\_v3.py}. The
core routine is \texttt{extract\_bio\_tags()}, supported by
\texttt{normalize\_token()} and \texttt{get\_entity\_spans()}. The pipeline
assigns every token one of three \textbf{BIO} labels:

\begin{center}
\begin{tabular}{@{}ll@{}}
\toprule
\textbf{Tag} & \textbf{Meaning} \\
\midrule
\texttt{B-PRODUCT} & First (or only) token of a product name \\
\texttt{I-PRODUCT} & Continuation token of a multi-word product name \\
\texttt{O}         & Non-entity token \\
\bottomrule
\end{tabular}
\end{center}

\noindent A single-word entity is a lone \texttt{B-PRODUCT}; a multi-word entity
is \texttt{B-PRODUCT} followed by one or more \texttt{I-PRODUCT} tokens.

\subsection{Tokenisation and tag recognition}
The tagged review is split on whitespace (\texttt{text.split()}), giving a linear
list of raw tokens. The routine scans them left to right with a single index:

\begin{itemize}
  \item If a token begins with the prefix \texttt{\_NE\_}, the prefix is removed
        (\texttt{entity\_word = token[4:]}) and the token is labelled
        \texttt{B-PRODUCT}, opening a new entity span.
  \item If a token is punctuation-only (checked by \texttt{is\_punctuation\_only()})
        or empty, it is discarded --- it receives no tag and is excluded from all
        counts.
  \item Any other token is labelled \texttt{O}.
\end{itemize}

\subsection{Continuation matching and the whitelist}
Immediately after a \texttt{B-PRODUCT} token is opened, an inner loop inspects
the following tokens to decide whether the entity extends. A candidate token is
merged as \texttt{I-PRODUCT} only if \emph{all} of the following hold:

\begin{enumerate}
  \item it does \emph{not} itself begin with \texttt{\_NE\_} (that would start a
        new entity);
  \item its normalised, lower-cased form is present in the fixed whitelist
        \texttt{PRODUCT\_SECOND\_TOKENS};
  \item it passes the three noise guards (Section~\ref{sec:guards}).
\end{enumerate}

The loop keeps consuming continuation tokens until one fails a test, at which
point the entity closes.

\begin{note}[The whitelist is the decisive component]
A continuation token is merged \emph{only} if it matches
\texttt{PRODUCT\_SECOND\_TOKENS}, a hand-curated set of product ``type'' words.
It contains both Bangla and English forms:

\medskip
\begin{tabular}{@{}p{2.0cm}p{11.5cm}@{}}
\textbf{Bangla} & \bn{সাবান}, \bn{ওয়াশ}, \bn{অয়েল}, \bn{ওয়েল}, \bn{লোশন},
  \bn{ক্রিম}, \bn{শ্যাম্পু}, \bn{টাচ}, \bn{মাস্ক}, \bn{ট্রিমার}, \bn{ব্রাশ},
  \bn{ব্যাগ}, \bn{তেল}, \bn{বার}, \bn{ফোম}, \bn{পাউডার}, \bn{জেল}, \bn{সিরাম},
  \bn{টোনার}, \bn{স্প্রে}, \bn{বাম}, \bn{বাটার}, \bn{ড্রপ}, \bn{ট্যাবলেট},
  \bn{ক্যাপসুল}, \bn{সিরাপ}, \bn{সোপ} \\[4pt]
\textbf{English} & oil, wash, cream, soap, mask, lotion, shampoo, powder, gel,
  spray, serum, toner, balm, butter, foam, bar, drop, tablet, capsule, syrup,
  trimmer, brush \\
\end{tabular}

\medskip
Any continuation word outside this set cannot form a multi-word entity, which
accounts for most of the gap in Section~\ref{sec:gap}.
\end{note}

\subsection{Token normalisation}
Before a token is tested against the whitelist (or stored as an entity stem) it
passes through \texttt{normalize\_token()}, which applies three steps in order:

\begin{enumerate}
  \item \textbf{Split on interior sentence boundaries.} The token is cut at the
        first occurrence of \bn{।} \texttt{. ? ! ;} \bn{৷} and only the part
        before it is kept, e.g.\ \bn{ওয়াশ।রিকমন্ডেড} $\rightarrow$ \bn{ওয়াশ}.
  \item \textbf{Strip trailing punctuation.} Trailing commas, colons, brackets
        and quotes are removed, e.g.\ \bn{ওয়াশ,} $\rightarrow$ \bn{ওয়াশ}.
  \item \textbf{Strip one inflectional suffix} (longest match first), reducing
        the surface form to its stem.
\end{enumerate}

The recognised Bangla suffixes are:

\begin{center}
\begin{tabular}{@{}lll@{}}
\toprule
\textbf{Suffix} & \textbf{Function} & \textbf{Example} \\
\midrule
\bn{গুলো} / \bn{গুলি} / \bn{গুলা} & plural classifier   & \bn{সাবানগুলা} $\rightarrow$ \bn{সাবান} \\
\bn{খানা} / \bn{খানি}             & unit classifier     & \bn{বইখানা} $\rightarrow$ \bn{বই} \\
\bn{টা} / \bn{টি} / \bn{টো}       & definiteness marker & \bn{ক্রিমটি} $\rightarrow$ \bn{ক্রিম} \\
\bn{ের}                            & genitive            & \bn{সাবানের} $\rightarrow$ \bn{সাবান} \\
\bn{কে}                            & dative / accusative & --- \\
\bn{তে}                            & locative            & --- \\
\bn{র}                             & bare genitive       & --- \\
\bn{ও}                             & additive            & --- \\
\bottomrule
\end{tabular}
\end{center}

\noindent Only one suffix is removed per token, and only when the remaining stem
is long enough to be meaningful (the code requires the stem to exceed the suffix
length). The genitive suffix uses the dependent vowel sign \bn{ে} (U+09C7)
rather than the independent vowel \bn{এ}; using the wrong codepoint would cause
silent match failures.

\subsection{Noise guards}
\label{sec:guards}
Three guards prevent spurious multi-word entities during the continuation loop:

\begin{center}
\begin{tabular}{@{}p{1.1cm}p{12.4cm}@{}}
\toprule
\textbf{G1} & \textbf{Same-token repetition.} A candidate whose stem is already
              present in the current span is rejected, so \bn{\_NE\_সাবান সাবান}
              yields a single entity \bn{সাবান} rather than the false pair
              \bn{সাবান সাবান}. \\[3pt]
\textbf{G2} & \textbf{Interior sentence boundary.} If the \texttt{B-PRODUCT}
              token itself contains a sentence-boundary mark in its interior
              (e.g.\ \bn{\_NE\_পণ্য।ইতিপূর্বেও ব্যাগ}), no continuation is
              attempted at all --- the merge would cross a sentence. \\[3pt]
\textbf{G3} & \textbf{Comma-list guard.} Using the \emph{original} untagged
              review, if the current stem is immediately followed by a comma
              (\bn{তেল, শ্যাম্পু}), the tokens are treated as an enumeration of
              separate products, not one entity. \\
\bottomrule
\end{tabular}
\end{center}

\subsection{Span extraction and deduplication}
After tagging, \texttt{get\_entity\_spans()} walks the tag sequence and groups
each \texttt{B-PRODUCT} with its trailing \texttt{I-PRODUCT} tokens into a span.
Each span stores its normalised token list, the joined surface text, its start
and end indices, and an \texttt{is\_multi} flag (true when it holds more than one
token). The joined \emph{normalised} text is the deduplication key: unique
multi-word entities are the distinct normalised strings, counted with a
\texttt{Counter}. Because the key is an exact string, two spellings of one
product (e.g.\ \bn{স্যাভলন সাবান} and \bn{Savlon সাবান}) form two separate keys.

\subsection{Worked example}
Input tag stream: \bn{\_NE\_ডেটল} \bn{সাবানটি} \bn{তো} \bn{ভাল}

\begin{enumerate}
  \item \bn{\_NE\_ডেটল} $\rightarrow$ prefix removed, stem \bn{ডেটল}, labelled
        \texttt{B-PRODUCT}; a new span opens.
  \item \bn{সাবানটি}: normalisation strips the definiteness suffix \bn{টি}
        $\rightarrow$ \bn{সাবান}; this is in the whitelist and passes G1--G3, so
        it is labelled \texttt{I-PRODUCT} and merged.
  \item \bn{তো}: not whitelisted $\rightarrow$ the span closes; token labelled
        \texttt{O}.
  \item \bn{ভাল}: labelled \texttt{O}.
\end{enumerate}

\noindent Resulting span: \bn{ডেটল সাবান} --- tokens \bn{ডেটল} + \bn{সাবান},
tags \texttt{B-PRODUCT + I-PRODUCT}, \texttt{is\_multi = true}.


% =============================================================================
\section{Results}

\begin{center}
\begin{tabular}{@{}lr@{}}
\toprule
\textbf{Measure} & \textbf{Value} \\
\midrule
Total reviews                    & 10,000 \\
Total tokens                     & 86,741 \\
\texttt{B-PRODUCT} tags          & 8,821 \\
\texttt{I-PRODUCT} tags          & 55 \\
\texttt{O} tags                  & 77,865 \\
Total entity spans               & 8,821 \\
\quad Single-word                & 8,766 \ (99.4\%) \\
\quad Multi-word                 & 55 \ \ \ (0.6\%) \\
Unique multi-word entities       & \textbf{13} \\
\bottomrule
\end{tabular}
\end{center}

\noindent The 13 unique multi-word entities detected:

\begin{center}
\small
\begin{tabular}{@{}rl r rl@{}}
\toprule
\textbf{\#} & \textbf{Entity} & \textbf{Count} & \textbf{\#} & \textbf{Entity} \\
\midrule
1 & \bn{ফেস ওয়াশ}      & 28 & 8  & \bn{প্যানটিন শ্যাম্পু} \\
2 & \bn{ডেটল সাবান}     & 14 & 9  & \bn{থ্রিডি মাস্ক} \\
3 & \bn{ফেস মাস্ক}      & 2  & 10 & \bn{সানছ ক্রিম} \\
4 & \bn{সেপনিল মাস্ক}   & 2  & 11 & \bn{ফেইস ওয়াশ} \\
5 & \bn{লাক্স সাবান}    & 1  & 12 & \bn{অলিভ অয়েল} \\
6 & \bn{স্যাভলন সাবান}  & 1  & 13 & \bn{ডেটল সোপ} \\
7 & \bn{Savlon সাবান}   & 1  &    & \\
\bottomrule
\end{tabular}
\end{center}

Three of these are spelling variants of products already in the list:
\bn{Savlon সাবান} duplicates \bn{স্যাভলন সাবান}, \bn{ফেইস ওয়াশ} duplicates
\bn{ফেস ওয়াশ}, and \bn{ডেটল সোপ} duplicates \bn{ডেটল সাবান}. The pipeline treats
each spelling as a distinct entity because it matches text exactly.

% =============================================================================
\section{Manual Review}

Reading the reviews directly, the human reviewer catalogued \textbf{30} unique
multi-word product names (\texttt{manual\_report.txt}):

\begin{center}
\small
\begin{tabular}{@{}rl rl rl@{}}
\toprule
\# & Entity & \# & Entity & \# & Entity \\
\midrule
1  & \bn{ফেস ওয়াশ}       & 11 & \bn{মেথি গুড়ো}     & 21 & \bn{টুথপেস্ট} \\
2  & \bn{ডেটল সাবান}      & 12 & \bn{আইলাইনার}       & 22 & \bn{টুথব্রাশ} \\
3  & \bn{ফেসমাস্ক}        & 13 & \bn{ফেসক্রিম}       & 23 & \bn{টি-শার্ট} \\
4  & \bn{সেপনিল মাস্ক}    & 14 & \bn{হেয়ার জেল}     & 24 & \bn{বডি স্প্রে} \\
5  & \bn{লাক্স সাবান}     & 15 & \bn{লিপজেল}         & 25 & \bn{শেভিং মেশিন} \\
6  & \bn{স্যাভলন সাবান}   & 16 & \bn{হোয়াইট বিউটি}  & 26 & \bn{মেহেদী গুড়ো} \\
7  & \bn{প্যানটিন শ্যাম্পু} & 17 & \bn{স্লিপিং মাস্ক} & 27 & \bn{সেপটেক্স সাবান} \\
8  & \bn{থ্রিডি মাস্ক}    & 18 & \bn{আলোভেরা জেল}    & 28 & \bn{মেরিলের সাবান} \\
9  & \bn{সানছ ক্রিম}      & 19 & \bn{কটন প্যাড}      & 29 & \bn{পেট্রোলিয়াম জেলি} \\
10 & \bn{অলিভ অয়েল}      & 20 & \bn{সার্জিক্যাল মাস্ক} & 30 & \bn{বেবি লোশন} \\
\bottomrule
\end{tabular}
\end{center}

% =============================================================================
\section{Gap Analysis: 13 vs.\ 30}
\label{sec:gap}

The pipeline captured 9 of the 30 products as exact multi-word matches. The
remaining 21 were missed for three reasons.

\subsection{Continuation word absent from the whitelist}
The pipeline merges a second token only if it is listed in
\texttt{PRODUCT\_SECOND\_TOKENS}. Products ending in unlisted words could not be
formed:

\begin{center}
\begin{tabular}{@{}ll@{}}
\toprule
\textbf{Missing continuation word} & \textbf{Affected products} \\
\midrule
\bn{গুড়ো} (powder)  & \bn{মেথি গুড়ো}, \bn{মেহেদী গুড়ো} \\
\bn{বিউটি} (beauty)  & \bn{হোয়াইট বিউটি} \\
\bn{প্যাড} (pad)     & \bn{কটন প্যাড} \\
\bn{মেশিন} (machine) & \bn{শেভিং মেশিন} \\
\bn{জেলি} (jelly)    & \bn{পেট্রোলিয়াম জেলি} \\
\bottomrule
\end{tabular}
\end{center}

\subsection{Product written as a single joined token}
Several products are written without a space and so are tokenised as one word.
The pipeline classifies them as single-word entities and never as multi-word:
\bn{ফেসমাস্ক} (face+mask), \bn{ফেসক্রিম} (face+cream), \bn{লিপজেল} (lip+gel),
\bn{আইলাইনার} (eye+liner), \bn{টুথপেস্ট} (tooth+paste), \bn{টুথব্রাশ}
(tooth+brush), \bn{টি-শার্ট} (t-shirt). It cannot split a token that contains no
space. Note that the code did detect the spaced form \bn{ফেস মাস্ক} elsewhere,
but not the joined \bn{ফেসমাস্ক} the reviewer recorded.

\subsection{Spelling and transliteration variants}
The same product appears under different spellings and scripts. Because matching
is exact, the pipeline either counts variants separately or fails to link them.
For example, the single product \textbf{Savlon Soap} occurs as
\bn{স্যাভলন সাবান} and \bn{Savlon সাবান}, and could equally appear as \bn{সালভন}
or \emph{Savlon Soap}. The reviewer recognises these as one product; the pipeline
does not, because it cannot equate \bn{সাবান} with ``soap'' or \bn{ফেস} with
\bn{ফেইস}. The same applies to \bn{ডেটল সাবান} vs.\ \bn{ডেটল সোপ}.

\subsection{Tagged but not merged}
A few products whose continuation word \emph{is} whitelisted (\bn{সাবান},
\bn{মাস্ক}, \bn{লোশন}, \bn{স্প্রে}, \bn{জেল}) still did not appear as merged
entities --- for example \bn{সেপটেক্স সাবান}, \bn{মেরিলের সাবান},
\bn{সার্জিক্যাল মাস্ক}, \bn{স্লিপিং মাস্ক}, \bn{বডি স্প্রে}, \bn{আলোভেরা জেল},
\bn{হেয়ার জেল}, \bn{বেবি লোশন}. In these cases the two tokens were not both
\texttt{\_NE\_}-tagged, were not adjacent, or the brand token was spelled
unusually.

\subsection{Per-entity outcome}

\begin{center}
\footnotesize
\begin{longtable}{@{}rlcl@{}}
\toprule
\textbf{\#} & \textbf{Product} & \textbf{Detected} & \textbf{Reason} \\
\midrule
\endhead
1  & \bn{ফেস ওয়াশ}         & \yes & whitelisted (\bn{ওয়াশ}) \\
2  & \bn{ডেটল সাবান}        & \yes & whitelisted (\bn{সাবান}) \\
3  & \bn{ফেসমাস্ক}          & \no  & written as one token \\
4  & \bn{সেপনিল মাস্ক}      & \yes & whitelisted (\bn{মাস্ক}) \\
5  & \bn{লাক্স সাবান}       & \yes & whitelisted (\bn{সাবান}) \\
6  & \bn{স্যাভলন সাবান}     & \yes & detected (also as \bn{Savlon সাবান}) \\
7  & \bn{প্যানটিন শ্যাম্পু} & \yes & whitelisted (\bn{শ্যাম্পু}) \\
8  & \bn{থ্রিডি মাস্ক}      & \yes & whitelisted (\bn{মাস্ক}) \\
9  & \bn{সানছ ক্রিম}        & \yes & whitelisted (\bn{ক্রিম}) \\
10 & \bn{অলিভ অয়েল}        & \yes & whitelisted (\bn{অয়েল}) \\
11 & \bn{মেথি গুড়ো}        & \no  & \bn{গুড়ো} not whitelisted \\
12 & \bn{আইলাইনার}          & \no  & written as one token \\
13 & \bn{ফেসক্রিম}          & \no  & written as one token \\
14 & \bn{হেয়ার জেল}        & \no  & not tagged / not merged \\
15 & \bn{লিপজেল}            & \no  & written as one token \\
16 & \bn{হোয়াইট বিউটি}     & \no  & \bn{বিউটি} not whitelisted \\
17 & \bn{স্লিপিং মাস্ক}     & \no  & not tagged / not merged \\
18 & \bn{আলোভেরা জেল}       & \no  & not tagged / not merged \\
19 & \bn{কটন প্যাড}         & \no  & \bn{প্যাড} not whitelisted \\
20 & \bn{সার্জিক্যাল মাস্ক} & \no  & not tagged / not merged \\
21 & \bn{টুথপেস্ট}          & \no  & written as one token \\
22 & \bn{টুথব্রাশ}          & \no  & written as one token \\
23 & \bn{টি-শার্ট}          & \no  & written as one token \\
24 & \bn{বডি স্প্রে}        & \no  & not tagged / not merged \\
25 & \bn{শেভিং মেশিন}       & \no  & \bn{মেশিন} not whitelisted \\
26 & \bn{মেহেদী গুড়ো}      & \no  & \bn{গুড়ো} not whitelisted \\
27 & \bn{সেপটেক্স সাবান}    & \no  & not tagged / not merged \\
28 & \bn{মেরিলের সাবান}     & \no  & not tagged / not merged \\
29 & \bn{পেট্রোলিয়াম জেলি} & \no  & \bn{জেলি} not whitelisted \\
30 & \bn{বেবি লোশন}         & \no  & not tagged / not merged \\
\bottomrule
\end{longtable}
\end{center}

% =============================================================================
\section{Limitations}

\subsection{Automated pipeline}
\begin{itemize}
  \item \textbf{Fixed whitelist.} Only continuation words in
        \texttt{PRODUCT\_SECOND\_TOKENS} can form a multi-word entity; unlisted
        words (\bn{গুড়ো}, \bn{বিউটি}, \bn{প্যাড}, \bn{মেশিন}, \bn{জেলি}) are missed.
  \item \textbf{Exact matching.} Spelling and transliteration variants
        (\bn{স্যাভলন} / Savlon / \bn{সালভন}) are not unified.
  \item \textbf{No token splitting.} Space-free compounds (\bn{ফেসমাস্ক},
        \bn{টুথপেস্ট}) are read as single words.
  \item \textbf{No cross-language equivalence.} \bn{সাবান} and ``soap'', or
        \bn{ফেস} and \bn{ফেইস}, are treated as unrelated.
  \item \textbf{Tag dependence.} Detection is impossible where the source
        \texttt{\_NE\_} tag is absent.
\end{itemize}

\subsection{Manual review}
\begin{itemize}
  \item \textbf{Not scalable} to large corpora; slow and labour-intensive.
  \item \textbf{Subjective and inconsistent} across reviewers or sessions.
  \item \textbf{Error-prone} under fatigue; occurrences can be miscounted or missed.
  \item \textbf{Not reproducible} and produces no automatic statistics
        (counts, sentiment breakdown, sentence-ID map).
\end{itemize}

% =============================================================================
\section{Conclusion}

The two methods measure different things and each is incomplete alone. The
pipeline is fast, consistent, and produces exact statistics, but is bounded by a
fixed whitelist, exact string matching, and dependence on source tags. Manual
review captures spelling variants, compounds, and unlisted product types through
human judgement, but does not scale and is not reproducible.

The practical improvement is to combine them: extend the whitelist (e.g.\
\bn{গুড়ো}, \bn{বিউটি}, \bn{প্যাড}, \bn{মেশিন}, \bn{জেলি}), add a
transliteration/normalisation layer so brand variants collapse to one entity,
add a compound-word splitter for space-free products, and retain a human review
pass over flagged edge cases using the sentence-ID map the pipeline already
produces.

\end{document}
